%% ============================================================================
%%  Colophon e licenza
%% ============================================================================
\clearpage
\thispagestyle{empty}
\null\vfill
\begin{flushleft}
\small
{\sffamily\bfseries\color{blu}Analisi Matematica} --- Un corso rigoroso, spiegato passo per passo.\\
Volume primo.\\[2mm]
Copyright \copyright{} \the\year{} Biagio Raucci.\\[4mm]

{\sffamily\bfseries\color{blu}Licenza.}
Quest'opera è rilasciata con licenza
\textbf{Creative Commons Attribuzione -- Condividi allo stesso modo 4.0 Internazionale}
(CC~BY-SA~4.0). Chiunque è libero di \emph{condividere} (riprodurre, distribuire,
comunicare al pubblico) e \emph{modificare} (remixare, trasformare il materiale e
basarsi su di esso per le proprie opere) per qualsiasi fine, anche commerciale, alle
seguenti condizioni: \emph{attribuzione} (è necessario riconoscere la paternità
dell'opera, fornire un link alla licenza e indicare se sono state effettuate delle
modifiche) e \emph{condivisione allo stesso modo} (se si trasforma il materiale, si deve
distribuire il proprio contributo con la stessa licenza dell'originale).
Il testo integrale della licenza è disponibile all'indirizzo\\
\url{https://creativecommons.org/licenses/by-sa/4.0/deed.it}\\[4mm]

{\sffamily\bfseries\color{blu}Sorgenti.}
Il libro è scritto in \LaTeX{} e i sorgenti --- testo, figure e codice Python ---
sono distribuiti insieme all'opera. La compilazione richiede \textsf{LuaLaTeX}:
\texttt{latexmk -lualatex main}.\\[4mm]

{\sffamily\bfseries\color{blu}Ringraziamenti e fonti.}
L'impostazione dei capitoli e la scelta degli argomenti seguono i classici corsi
universitari di Analisi Matematica Uno; in particolare i testi di P.~Marcellini e
C.~Sbordone, di E.~Giusti e di T.~Tao, e le dispense \emph{Analisi Uno} di E.~Paolini
(anch'esse rilasciate con licenza CC~BY-SA), che hanno ispirato la struttura di questo
progetto. Gli ornamenti floreali delle pagine di apertura sono generati dal pacchetto
\textsf{florilegium} di Kevin Raymond (licenza MIT), che disegna ogni petalo e ogni
foglia in Lua al momento della composizione: nessuna immagine è importata.\\[4mm]

{\sffamily\bfseries\color{blu}Software.}
Composto con \LaTeX{} (motore LuaTeX), font TeX Gyre Pagella e TeX Gyre Heros.
Figure realizzate con Ti\emph{k}Z, PGFPlots e Matplotlib.
\end{flushleft}
\vspace{10mm}
